\documentclass[CJK,13pt]{beamer}
\usepackage{CJKutf8}
\usepackage{beamerthemesplit}
\usetheme{Malmoe}
\useoutertheme[footline=authortitle]{miniframes}
\usepackage{amsmath}
\usepackage{amssymb}
\usepackage{graphicx}
\usepackage{eufrak}
\usepackage{color}
\usepackage{slashed}
\usepackage{simplewick}
\usepackage{tikz}
\usepackage{tcolorbox}
\usepackage{ulem}
\graphicspath{{../figures/}}
%%figures
\def\lfig#1#2{\includegraphics[width=#1 in]{#2}}
\def\tfig#1#2{\includegraphics[height=#1 in]{#2}}
\def\addfig#1#2{\begin{center}\includegraphics[width=#1 in]{#2}\end{center}}
\def\question{\includegraphics[width=0.3in]{why.jpg}\,}
\def\answer{\includegraphics[width=0.3in]{baozou_haha.png}\,}
\def\wulian{\includegraphics[width=0.18in]{emoji_wulian.jpg}}
\def\bigwulian{\includegraphics[width=0.35in]{emoji_wulian.jpg}}
\def\bye{\includegraphics[width=0.18in]{emoji_bye.jpg}}
\def\bigbye{\includegraphics[width=0.35in]{emoji_bye.jpg}}
\def\huaixiao{\includegraphics[width=0.18in]{emoji_huaixiao.jpg}}
\def\bighuaixiao{\includegraphics[width=0.35in]{emoji_huaixiao.jpg}}
\def\jianxiao{\includegraphics[width=0.18in]{emoji_jianxiao.jpg}}
\def\bigjianxiao{\includegraphics[width=0.35in]{emoji_jianxiao.jpg}}
\def\haoqi{\includegraphics[width=0.18in]{emoji_haoqi.jpg}}
%% colors
\def\blacktext#1{{\color{black}#1}}
\def\bluetext#1{{\color{blue}#1}}
\def\redtext#1{{\color{red}#1}}
\def\darkbluetext#1{{\color[rgb]{0,0.2,0.6}#1}}
\def\skybluetext#1{{\color[rgb]{0.2,0.7,1.}#1}}
\def\cyantext#1{{\color[rgb]{0.,0.5,0.5}#1}}
\def\greentext#1{{\color[rgb]{0,0.7,0.1}#1}}
\def\darkgray{\color[rgb]{0.2,0.2,0.2}}
\def\lightgray{\color[rgb]{0.6,0.6,0.6}}
\def\gray{\color[rgb]{0.4,0.4,0.4}}
\def\blue{\color{blue}}
\def\red{\color{red}}
\def\orange{\color[rgb]{1.,0.8,0.}}
\def\green{\color{green}}
\def\darkgreen{\color[rgb]{0,0.4,0.1}}
\def\darkblue{\color[rgb]{0,0.2,0.6}}
\def\skyblue{\color[rgb]{0.2,0.7,1.}}
%%control
\def\diag{\mathrm{diag}\,}
\def\heaviside{\,\mathrm{h}\,}
\def\bral{\left(\begin{array}{l}}
\def\brar{\end{array}\right)}
\def\brall{\left(\begin{array}{ll}}
\def\brarr{\end{array}\right)}
\def\bralll{\left(\begin{array}{lll}}
\def\brarrr{\end{array}\right)}
\def\branchl{\left\{\begin{array}{l}}
\def\branchr{\end{array}\right.}
\def\branchll{\left\{\begin{array}{ll}}
\def\branchrr{\end{array}\right.}
\def\branchlll{\left\{\begin{array}{lll}}
\def\branchrrr{\end{array}\right.}
\def\sfgamma#1{\,\Gamma\left( #1 \right)\,}
\def\be{\begin{equation}}
\def\ee{\nonumber\end{equation}}
\def\bea{\begin{eqnarray}}
\def\eea{\nonumber\end{eqnarray}}
\def\bch{\begin{CJK}{UTF8}{gbsn}}
\def\ech{\end{CJK}}
\def\bitem{\begin{itemize}}
\def\eitem{\end{itemize}}
\def\bcenter{\begin{center}}
\def\ecenter{\end{center}}
\def\bex{\begin{minipage}{0.2\textwidth}\includegraphics[width=0.6in]{jugelizi.png}\end{minipage}\begin{minipage}{0.76\textwidth}}
\def\eex{\end{minipage}}
\def\chtitle#1{\frametitle{\bch#1\ech}}
\def\bmat#1{\left(\begin{array}{#1}}
\def\emat{\end{array}\right)}
\def\bcase#1{\left\{\begin{array}{#1}}
\def\ecase{\end{array}\right.}
\def\bmini#1{\begin{minipage}{#1\textwidth}}
\def\emini{\end{minipage}}
\def\tbox#1{\begin{tcolorbox}#1\end{tcolorbox}}
\def\pfrac#1#2#3{\left(\frac{\partial #1}{\partial #2}\right)_{#3}}
\def\res#1#2{\,\mathrm{res}\,#1\left(#2\right)\,}
\def\newt#1#2{\left(\begin{array}{c}#1\\ #2\end{array}\right)}
\def\reof#1{\,\mathrm{Re}{\left(#1\right)\,}}
\def\imof#1{\,\mathrm{Im}{\left(#1\right)\,}}
\def\Arg#1{\,\mathrm{Arg}\,#1\,}
%%symbols
\def\intfull{\int_{-\infty}^\infty}
\def\inthalf{\int_0^\infty}
\def\bropt{\,(\ \ \ )}
\def\sech{\mathrm{sech}\,}
\def\csch{\mathrm{csch}\,}
\def\asinh{\mathrm{asinh}\,}
\def\acosh{\mathrm{acosh}\,}
\def\sone{$\star$}
\def\stwo{$\star\star$}
\def\sthree{$\star\star\star$}
\def\sfour{$\star\star\star\star$}
\def\sfive{$\star\star\star\star\star$}
\def\rint{{\int_\leftrightarrow}}
\def\roint{{\oint_\leftrightarrow}}
\def\stdHf{{\textit{\r H}_f}}
\def\deltaH{{\Delta \textit{\r H}}}
\def\ii{{\dot{\imath}}}
\def\skipline{{\vskip0.1in}}
\def\skiplines{{\vskip0.2in}}
\def\lagr{{\mathcal{L}}}
\def\hamil{{\mathcal{H}}}
\def\vecT{{\mathbf{T}}}
\def\vecN{{\mathbf{N}}}
\def\vecB{{\mathbf{B}}}
\def\vecv{{\mathbf{v}}}
\def\vecr{{\mathbf{r}}}
\def\vecf{{\mathbf{f}}}
\def\vecg{{\mathbf{g}}}
\def\vecupsilon{{\mathbf{\upsilon}}}
\def\vecu{{\mathbf{u}}}
\def\vecj{{\mathbf{j}}}
\def\vecx{{\mathbf{x}}}
\def\vecy{{\mathbf{y}}}
\def\vecz{{\mathbf{z}}}
\def\veck{{\mathbf{k}}}
\def\vecp{{\mathbf{p}}}
\def\vecn{{\mathbf{n}}}
\def\vecA{{\mathbf{A}}}
\def\vecP{{\mathbf{P}}}
\def\vecsigma{{\mathbf{\sigma}}}
\def\hatJn{{\hat{J_\vecn}}}
\def\hatJx{{\hat{J_x}}}
\def\hatJy{{\hat{J_y}}}
\def\hatJz{{\hat{J_z}}}
\def\hatj#1{\hat{J_{#1}}}
\def\hatphi{{\hat{\phi}}}
\def\hatq{{\hat{q}}}
\def\hatpi{{\hat{\pi}}}
\def\vel{\upsilon}
\def\Dint{{\mathcal{D}}}
\def\adag{{\hat{a}^\dagger}}
\def\bdag{{\hat{b}^\dagger}}
\def\cdag{{\hat{c}^\dagger}}
\def\ddag{{\hat{d}^\dagger}}
\def\hata{{\hat{a}}}
\def\hatb{{\hat{b}}}
\def\hatc{{\hat{c}}}
\def\hatd{{\hat{d}}}
\def\hatD{{\,\hat{D}}}
\def\hatN{{\hat{N}}}
\def\hatH{{\hat{H}}}
\def\hatp{{\hat{p}}}
\def\Fup{{F^{\mu\nu}}}
\def\Fdown{{F_{\mu\nu}}}
\def\newl{\nonumber \\}
\def\vece{\mathrm{e}}
\def\calM{{\mathcal{M}}}
\def\calT{{\mathcal{T}}}
\def\calR{{\mathcal{R}}}
\def\barpsi{\bar{\psi}}
\def\baru{\bar{u}}
\def\barv{\bar{\upsilon}}
\def\qeq{\stackrel{?}{=}}
\def\ftf{\stackrel{\mathcal{FT}}{\Longrightarrow}}
\def\ftb{\stackrel{\mathcal{FT}}{\Longleftarrow}}
\def\ftfb{\stackrel{\mathcal{FT}}{\Longleftrightarrow}}
\def\ltf{\stackrel{\mathcal{LT}}{\Longrightarrow}}
\def\ltb{\stackrel{\mathcal{LT}}{\Longleftarrow}}
\def\ltfb{\stackrel{\mathcal{LT}}{\Longleftrightarrow}}
\def\torder#1{\mathcal{T}\left(#1\right)}
\def\rorder#1{\mathcal{R}\left(#1\right)}
\def\contr#1#2{\contraction{}{#1}{}{#2}#1#2}
\def\trof#1{\mathrm{Tr}\left(#1\right)}
\def\trace{\mathrm{Tr}}
\def\comm#1{\ \ \ \left(\mathrm{used}\ #1\right)}
\def\tcomm#1{\ \ \ (\text{#1})}
\def\slp{\slashed{p}}
\def\slk{\slashed{k}}
\def\calp{{\mathfrak{p}}}
\def\veccalp{\mathbf{\mathfrak{p}}}
\def\Tthree{T_{\tiny \textcircled{3}}}
\def\pthree{p_{\tiny \textcircled{3}}}
\def\dbar{{\,\mathchar'26\mkern-12mu d}}
\def\erf{\mathrm{erf}}
\def\const{\mathrm{const.}}
\def\pheat{\pfrac p{\ln T}V}
\def\vheat{\pfrac V{\ln T}p}
%%units
\def\fdeg{{^\circ \mathrm{F}}}
\def\cdeg{^\circ \mathrm{C}}
\def\atm{\,\mathrm{atm}}
\def\angstrom{\,\text{\AA}}
\def\SIL{\,\mathrm{L}}
\def\SIkm{\,\mathrm{km}}
\def\SIyr{\,\mathrm{yr}}
\def\SIGyr{\,\mathrm{Gyr}}
\def\SIV{\,\mathrm{V}}
\def\SImV{\,\mathrm{mV}}
\def\SIeV{\,\mathrm{eV}}
\def\SIkeV{\,\mathrm{keV}}
\def\SIMeV{\,\mathrm{MeV}}
\def\SIGeV{\,\mathrm{GeV}}
\def\SIcal{\,\mathrm{cal}}
\def\SIkcal{\,\mathrm{kcal}}
\def\SImol{\,\mathrm{mol}}
\def\SIN{\,\mathrm{N}}
\def\SIHz{\,\mathrm{Hz}}
\def\SIm{\,\mathrm{m}}
\def\SIcm{\,\mathrm{cm}}
\def\SIfm{\,\mathrm{fm}}
\def\SImm{\,\mathrm{mm}}
\def\SInm{\,\mathrm{nm}}
\def\SImum{\,\mathrm{\mu m}}
\def\SIJ{\,\mathrm{J}}
\def\SIW{\,\mathrm{W}}
\def\SIkJ{\,\mathrm{kJ}}
\def\SIs{\,\mathrm{s}}
\def\SIkg{\,\mathrm{kg}}
\def\SIg{\,\mathrm{g}}
\def\SIK{\,\mathrm{K}}
\def\SImmHg{\,\mathrm{mmHg}}
\def\SIPa{\,\mathrm{Pa}}
%page
\def\secpage#1#2{\begin{frame}\bcenter{\bf \Huge #1} \ecenter \skipline \tbox{\bcenter #2 \ecenter}\end{frame}}
\def\append#1#2{\secpage{附录#1}{#2}}
\def\thinka#1{\begin{frame}\frametitle{思考题}\bcenter\lfig{0.4}{think0.jpg}  \skipline #1 \ecenter\end{frame}}
\def\thinkb#1{\begin{frame}\frametitle{思考题}\bcenter\lfig{0.5}{think1.jpg}  \skipline #1 \ecenter\end{frame}}
\def\thinkc#1{\begin{frame}\frametitle{思考题}\bcenter\lfig{0.5}{think2.jpg}  \skipline #1 \ecenter\end{frame}}
\def\thinkd#1{\begin{frame}\frametitle{思考题}\bcenter\lfig{0.5}{think3.jpg}  \skipline #1 \ecenter\end{frame}}
\def\thinke#1{\begin{frame}\frametitle{思考题}\bcenter\lfig{0.5}{think4.jpg}  \skipline #1 \ecenter\end{frame}}
\def\thinkf#1{\begin{frame}\frametitle{思考题}\bcenter\lfig{0.5}{think5.jpg}  \skipline #1 \ecenter\end{frame}}
\def\schw{ds^2 = \left(1-\frac{2GM}{r}\right)dt^2 - \left(1-\frac{2GM}{r}\right)^{-1} dr^2 - r^2\left(d\theta^2 + \sin^2\theta d\phi^2\right)}

\input{titlepage.tex}
  \date{}
  \begin{document}
  \bch
  \tpage{20}{Gravitational Waves}


  \begin{frame}
    往水里扔一块石头，水的流体元的空间位置随着时间波动性地变化。

    \addfig{1.5}{waterwave.jpg}

    可以这样说：水波是 \uline{水} 相对于 \uline{空间背景} 的 \uline{时间}周期性 振动。
  \end{frame}


  \begin{frame}
    \frametitle{问题I}
    引力波通常在科普读物中被描述为“时空的涟漪”
    
    \addfig{1.5}{yinlibo.jpg}

    那引力波是 \uline{？？} 相对于 \uline{？？？} 的 \uline{？？} 周期性 振动？
  \end{frame}


  \begin{frame}
    \frametitle{问题II}
    爱因斯坦方程
    $$R^{\mu\nu}-\frac{1}{2}Rg^{\mu\nu} = 8\pi G T^{\mu\nu}$$
    这里的 $T^{\mu\nu}$ 要包含引力波的能量动量吗？
  \end{frame}


  \begin{frame}
    \frametitle{问题III}
    引力波经过引力源附近也会发生偏折吗？
  \end{frame}


  \begin{frame}
    谁说知道引力波是啥，你就拿这三连发锤TA。

    \addfig{0.8}{lahei.jpg}
  \end{frame}

  \secpage{谐和坐标系的引力波方程}{$$\square h_{\mu\nu} = -16\pi G\left(T_{\mu\nu}- \frac{1}{2}\eta_{\mu\nu}T\right)$$}

    \begin{frame}
    \frametitle{四维拉普拉斯算符$\square$}
    四维拉普拉斯算符 $\square$ 的定义为
    $$\square f = f^{;\mu}_{\ \ ;\mu}$$
    对一个标量 $\phi$，有
    \bea
    \square \phi &=& \left(g^{\mu\nu}\phi_{;\nu}\right)_{;\mu} \newl
    &=& g^{\mu\nu}\phi_{;\nu;\mu} \newl
    &=& g^{\mu\nu}\phi_{,\nu,\mu} - g^{\mu\nu}\Gamma^\lambda_{\ \nu\mu}\phi_{,\lambda} \newl
    &=& g^{\mu\nu}\phi_{,\nu,\mu} - \Gamma^\lambda\phi_{,\lambda} 
    \eea
    在最后一行我们定义了 $\Gamma^\lambda\equiv g^{\mu\nu}\Gamma^\lambda_{\mu\nu}$
  \end{frame}


    \begin{frame}

      
      显然，Minkowski时空的里相速度为光速的波动方程可以写成 $\square u = 0$.

      \addfig{1}{haowubodong.jpg}


      
  \end{frame}
    
  
  \begin{frame}
    \frametitle{近似Minkowski时空}
    如果观测者远离中子星和黑洞这样稀少的天体。我们知道广义相对论效应都不强，时空近似是 Minkowski 时空。

    \skipline
    
    不仅如此，爱因斯坦方程右边的 $8\pi G T^{\mu\nu}$ 的量级可以用白矮星的密度来设置一个上限：
    $$\lvert 8\pi G T^{\mu\nu}\rvert \lesssim \frac{1}{L^2}$$
    这里的 $L=10^9\mathrm{m}$。由于爱因斯坦张量是度规的二阶导数以及一阶导数的平方组成的，所以就可以认为在小于 $L$ 的尺度范围内度规的变化量也很小。这当然不能算是严格的证明，但至少是作为推理出发点的合理假设。
  \end{frame}


  \begin{frame}
    \frametitle{近Minkowski度规}
    因此，我们假设在远离致密天体的小范围时空区域里，{\bf 选择合理的坐标系}，可以有度规
    $$ g_{\mu\nu} = \eta_{\mu\nu}+h_{\mu\nu}$$
    这里的 $\eta_{\mu\nu}$ 是Minkowski度规，$h_{\mu\nu}\ll 1$，且 $h_{\mu\nu}$ 的一阶、二阶偏导数都可以看成同阶的小量。我们将在讨论运动方程时总是{\blue 只保留 $h_{\mu\nu}$ 的一阶小量}。
  \end{frame}


  \begin{frame}
    \frametitle{“合理坐标系”不清不楚，不如叫“\sout{河蟹}谐和坐标系”}
    我们要求“谐和坐标系”对所有指标$\mu=0,1,2,3$，都要满足
  
    \addfig{0.8}{crab.jpg}
  
    取标量 $f = x^\mu$ （那就是说，即使你换到另外一种不太河蟹的坐标系，每个物理点的 $f$ 还是原先谐和坐标系里的 $x^\mu$），都要有
    $$\square f = 0$$
  \end{frame}

  \begin{frame}
    \frametitle{谐和坐标条件（续）}
    根据
    $$ \square f = g^{\mu\nu}f_{,\mu,\nu} - \Gamma^\lambda f_{,\lambda}$$
    在谐和坐标系里， $f=x^\mu$ 的任意二阶普通偏导都是零，所以谐和坐标条件可以等价地写为
    $$\Gamma^\lambda = 0$$
    当然，也可以写成
    $$\Gamma_\lambda \equiv g^{\alpha\beta}\Gamma_{\lambda\alpha\beta}=0$$
  \end{frame}




  \begin{frame}
    \frametitle{线性化的谐和坐标条件}
    我们先计算联络
    $$\Gamma_{\rho\mu\nu} = \frac{1}{2}\left(h_{\mu\rho,\nu}+h_{\nu\rho,\mu}-h_{\mu\nu,\rho}\right)$$    
    谐和坐标条件就是
    $$\frac{1}{2}\eta^{\mu\nu}\left(h_{\mu\rho,\nu}+h_{\nu\rho,\mu}-h_{\mu\nu,\rho}\right) = 0$$
    上式可以写成
    $$ \eta^{\mu\nu}h_{\rho\mu,\nu} =\frac{1}{2} h_{,\rho} $$
    这里的 $h\equiv \eta^{\mu\nu}h_{\mu\nu}$.
  \end{frame}


  \begin{frame}
    \frametitle{谐和坐标系里的Ricci张量一阶近似}
    由于联络已经是 $h_{\mu\nu}$ 的一阶小量，在黎曼张量的联络表达式里就可以扔掉联络的乘积项，
    {\small
    \bea
    R_{\mu\nu} &=& R^\lambda_{\ \mu\nu\lambda} \newl
    &=& \Gamma^\lambda_{\ \mu\nu,\lambda} - \Gamma^\lambda_{\ \mu\lambda,\nu} \newl
    &=&  \frac{1}{2}\eta^{\lambda\rho}\left(h_{\mu\rho,\nu,\lambda}+h_{\nu\rho,\mu,\lambda}-h_{\mu\nu,\rho,\lambda}+h_{\rho\mu,\lambda,\nu}+h_{\rho\lambda,\mu,\nu}-h_{\mu\lambda,\rho,\nu}\right) \newl
    &=&  \frac{1}{2}\eta^{\lambda\rho}\left(h_{\nu\rho,\mu,\lambda}+h_{\rho\mu,\lambda,\nu}-h_{\mu\nu,\rho,\lambda}\right) + \frac{1}{2}h_{,\mu,\nu}     \newl
    &=&  -\frac{1}{2}\eta^{\lambda\rho} h_{\mu\nu,\rho,\lambda}     \newl
    \eea
    }
    在最后一步我们使用了谐和坐标条件        
  \end{frame}


  \begin{frame}
    \frametitle{真空中的爱因斯坦方程：引力波}
    于是在谐和坐标系下，真空中的爱因斯坦方程可以写作：
    $$ \square h_{\mu\nu} = 0$$
    这是波速为光速的波动方程。
  \end{frame}


  \begin{frame}
    \frametitle{非真空情况}
    在非真空中，谐和坐标系里线性化的爱因斯坦方程可以写作：{\blue
      $$\square h_{\mu\nu} = S_{\mu\nu}$$}
    这里的
    $$ S_{\mu\nu} \equiv -16\pi G\left(T_{\mu\nu} - \frac{1}{2}\eta_{\mu\nu}T^\alpha_{\ \alpha}\right)$$
    为了让解看起来更幼儿园化，下面我们把时间和空间坐标分离开来，记 $t=x^0$, $\mathbf{x} = (x^1, x^2, x^3)$。    
  \end{frame}

  \begin{frame}
    \frametitle{引力辐射方程的解}
    如果你\uline{贫瘠的记忆里还有一丝电动力学的残余物}，或许下面的“推迟势”解并不会让你太慌张：
    $$ h_{\mu\nu}(t, \vecx) = \frac{1}{4\pi}\int d^3\vecx' \frac{S_{\mu\nu}\left(\vecx', t - \lvert \vecx -\vecx'\rvert\right)}{\lvert \vecx - \vecx'\rvert} $$

    \addfig{0.8}{wangji.jpg}
    
    当然这是一个特解，你可以把任何真空中的 $\square h_{\mu\nu}=0$ 的解加上去。和电动力学中一样，我们把这个特解解释为由源 $S_{\mu\nu}$ 产生的引力波，而把额外的真空解部分（如果有的话）解释为来自于远处未知源的引力波。
  \end{frame}
  
    \ech
\end{document}




  
